% Options for packages loaded elsewhere
\PassOptionsToPackage{unicode}{hyperref}
\PassOptionsToPackage{hyphens}{url}
%
\documentclass[
  letterpaper,
]{book}

\usepackage{amsmath,amssymb}
\usepackage{iftex}
\ifPDFTeX
  \usepackage[T1]{fontenc}
  \usepackage[utf8]{inputenc}
  \usepackage{textcomp} % provide euro and other symbols
\else % if luatex or xetex
  \usepackage{unicode-math}
  \defaultfontfeatures{Scale=MatchLowercase}
  \defaultfontfeatures[\rmfamily]{Ligatures=TeX,Scale=1}
\fi
\usepackage{lmodern}
\ifPDFTeX\else  
    % xetex/luatex font selection
\fi
% Use upquote if available, for straight quotes in verbatim environments
\IfFileExists{upquote.sty}{\usepackage{upquote}}{}
\IfFileExists{microtype.sty}{% use microtype if available
  \usepackage[]{microtype}
  \UseMicrotypeSet[protrusion]{basicmath} % disable protrusion for tt fonts
}{}
\makeatletter
\@ifundefined{KOMAClassName}{% if non-KOMA class
  \IfFileExists{parskip.sty}{%
    \usepackage{parskip}
  }{% else
    \setlength{\parindent}{0pt}
    \setlength{\parskip}{6pt plus 2pt minus 1pt}}
}{% if KOMA class
  \KOMAoptions{parskip=half}}
\makeatother
\usepackage{xcolor}
\usepackage[margin=1in]{geometry}
\setlength{\emergencystretch}{3em} % prevent overfull lines
\setcounter{secnumdepth}{5}
% Make \paragraph and \subparagraph free-standing
\makeatletter
\ifx\paragraph\undefined\else
  \let\oldparagraph\paragraph
  \renewcommand{\paragraph}{
    \@ifstar
      \xxxParagraphStar
      \xxxParagraphNoStar
  }
  \newcommand{\xxxParagraphStar}[1]{\oldparagraph*{#1}\mbox{}}
  \newcommand{\xxxParagraphNoStar}[1]{\oldparagraph{#1}\mbox{}}
\fi
\ifx\subparagraph\undefined\else
  \let\oldsubparagraph\subparagraph
  \renewcommand{\subparagraph}{
    \@ifstar
      \xxxSubParagraphStar
      \xxxSubParagraphNoStar
  }
  \newcommand{\xxxSubParagraphStar}[1]{\oldsubparagraph*{#1}\mbox{}}
  \newcommand{\xxxSubParagraphNoStar}[1]{\oldsubparagraph{#1}\mbox{}}
\fi
\makeatother

\usepackage{color}
\usepackage{fancyvrb}
\newcommand{\VerbBar}{|}
\newcommand{\VERB}{\Verb[commandchars=\\\{\}]}
\DefineVerbatimEnvironment{Highlighting}{Verbatim}{commandchars=\\\{\}}
% Add ',fontsize=\small' for more characters per line
\usepackage{framed}
\definecolor{shadecolor}{RGB}{241,243,245}
\newenvironment{Shaded}{\begin{snugshade}}{\end{snugshade}}
\newcommand{\AlertTok}[1]{\textcolor[rgb]{0.68,0.00,0.00}{#1}}
\newcommand{\AnnotationTok}[1]{\textcolor[rgb]{0.37,0.37,0.37}{#1}}
\newcommand{\AttributeTok}[1]{\textcolor[rgb]{0.40,0.45,0.13}{#1}}
\newcommand{\BaseNTok}[1]{\textcolor[rgb]{0.68,0.00,0.00}{#1}}
\newcommand{\BuiltInTok}[1]{\textcolor[rgb]{0.00,0.23,0.31}{#1}}
\newcommand{\CharTok}[1]{\textcolor[rgb]{0.13,0.47,0.30}{#1}}
\newcommand{\CommentTok}[1]{\textcolor[rgb]{0.37,0.37,0.37}{#1}}
\newcommand{\CommentVarTok}[1]{\textcolor[rgb]{0.37,0.37,0.37}{\textit{#1}}}
\newcommand{\ConstantTok}[1]{\textcolor[rgb]{0.56,0.35,0.01}{#1}}
\newcommand{\ControlFlowTok}[1]{\textcolor[rgb]{0.00,0.23,0.31}{\textbf{#1}}}
\newcommand{\DataTypeTok}[1]{\textcolor[rgb]{0.68,0.00,0.00}{#1}}
\newcommand{\DecValTok}[1]{\textcolor[rgb]{0.68,0.00,0.00}{#1}}
\newcommand{\DocumentationTok}[1]{\textcolor[rgb]{0.37,0.37,0.37}{\textit{#1}}}
\newcommand{\ErrorTok}[1]{\textcolor[rgb]{0.68,0.00,0.00}{#1}}
\newcommand{\ExtensionTok}[1]{\textcolor[rgb]{0.00,0.23,0.31}{#1}}
\newcommand{\FloatTok}[1]{\textcolor[rgb]{0.68,0.00,0.00}{#1}}
\newcommand{\FunctionTok}[1]{\textcolor[rgb]{0.28,0.35,0.67}{#1}}
\newcommand{\ImportTok}[1]{\textcolor[rgb]{0.00,0.46,0.62}{#1}}
\newcommand{\InformationTok}[1]{\textcolor[rgb]{0.37,0.37,0.37}{#1}}
\newcommand{\KeywordTok}[1]{\textcolor[rgb]{0.00,0.23,0.31}{\textbf{#1}}}
\newcommand{\NormalTok}[1]{\textcolor[rgb]{0.00,0.23,0.31}{#1}}
\newcommand{\OperatorTok}[1]{\textcolor[rgb]{0.37,0.37,0.37}{#1}}
\newcommand{\OtherTok}[1]{\textcolor[rgb]{0.00,0.23,0.31}{#1}}
\newcommand{\PreprocessorTok}[1]{\textcolor[rgb]{0.68,0.00,0.00}{#1}}
\newcommand{\RegionMarkerTok}[1]{\textcolor[rgb]{0.00,0.23,0.31}{#1}}
\newcommand{\SpecialCharTok}[1]{\textcolor[rgb]{0.37,0.37,0.37}{#1}}
\newcommand{\SpecialStringTok}[1]{\textcolor[rgb]{0.13,0.47,0.30}{#1}}
\newcommand{\StringTok}[1]{\textcolor[rgb]{0.13,0.47,0.30}{#1}}
\newcommand{\VariableTok}[1]{\textcolor[rgb]{0.07,0.07,0.07}{#1}}
\newcommand{\VerbatimStringTok}[1]{\textcolor[rgb]{0.13,0.47,0.30}{#1}}
\newcommand{\WarningTok}[1]{\textcolor[rgb]{0.37,0.37,0.37}{\textit{#1}}}

\providecommand{\tightlist}{%
  \setlength{\itemsep}{0pt}\setlength{\parskip}{0pt}}\usepackage{longtable,booktabs,array}
\usepackage{calc} % for calculating minipage widths
% Correct order of tables after \paragraph or \subparagraph
\usepackage{etoolbox}
\makeatletter
\patchcmd\longtable{\par}{\if@noskipsec\mbox{}\fi\par}{}{}
\makeatother
% Allow footnotes in longtable head/foot
\IfFileExists{footnotehyper.sty}{\usepackage{footnotehyper}}{\usepackage{footnote}}
\makesavenoteenv{longtable}
\usepackage{graphicx}
\makeatletter
\def\maxwidth{\ifdim\Gin@nat@width>\linewidth\linewidth\else\Gin@nat@width\fi}
\def\maxheight{\ifdim\Gin@nat@height>\textheight\textheight\else\Gin@nat@height\fi}
\makeatother
% Scale images if necessary, so that they will not overflow the page
% margins by default, and it is still possible to overwrite the defaults
% using explicit options in \includegraphics[width, height, ...]{}
\setkeys{Gin}{width=\maxwidth,height=\maxheight,keepaspectratio}
% Set default figure placement to htbp
\makeatletter
\def\fps@figure{htbp}
\makeatother

\setcounter{chapter}{0}
\renewcommand{\chaptername}{}
\makeatletter
\@ifpackageloaded{bookmark}{}{\usepackage{bookmark}}
\makeatother
\makeatletter
\@ifpackageloaded{caption}{}{\usepackage{caption}}
\AtBeginDocument{%
\ifdefined\contentsname
  \renewcommand*\contentsname{Table of contents}
\else
  \newcommand\contentsname{Table of contents}
\fi
\ifdefined\listfigurename
  \renewcommand*\listfigurename{List of Figures}
\else
  \newcommand\listfigurename{List of Figures}
\fi
\ifdefined\listtablename
  \renewcommand*\listtablename{List of Tables}
\else
  \newcommand\listtablename{List of Tables}
\fi
\ifdefined\figurename
  \renewcommand*\figurename{Figure}
\else
  \newcommand\figurename{Figure}
\fi
\ifdefined\tablename
  \renewcommand*\tablename{Table}
\else
  \newcommand\tablename{Table}
\fi
}
\@ifpackageloaded{float}{}{\usepackage{float}}
\floatstyle{ruled}
\@ifundefined{c@chapter}{\newfloat{codelisting}{h}{lop}}{\newfloat{codelisting}{h}{lop}[chapter]}
\floatname{codelisting}{Listing}
\newcommand*\listoflistings{\listof{codelisting}{List of Listings}}
\makeatother
\makeatletter
\makeatother
\makeatletter
\@ifpackageloaded{caption}{}{\usepackage{caption}}
\@ifpackageloaded{subcaption}{}{\usepackage{subcaption}}
\makeatother

\ifLuaTeX
  \usepackage{selnolig}  % disable illegal ligatures
\fi
\usepackage{bookmark}

\IfFileExists{xurl.sty}{\usepackage{xurl}}{} % add URL line breaks if available
\urlstyle{same} % disable monospaced font for URLs
\hypersetup{
  pdftitle={Practical One},
  hidelinks,
  pdfcreator={LaTeX via pandoc}}


\title{Practical One}
\author{}
\date{}

\begin{document}
\frontmatter
\maketitle

\renewcommand*\contentsname{Table of contents}
{
\setcounter{tocdepth}{1}
\tableofcontents
}

\mainmatter
\bookmarksetup{startatroot}

\chapter{Practical One}\label{practical-one}

\bookmarksetup{startatroot}

\chapter{}\label{section}

\chapter*{Question 1}

Find all rows in \textbf{``airquality''} data set that have missing
values.

\textbf{Code}

\begin{Shaded}
\begin{Highlighting}[]
\NormalTok{rows\_missing\_val }\OtherTok{\textless{}{-}}\NormalTok{airquality[}\FunctionTok{rowSums}\NormalTok{(}\FunctionTok{is.na}\NormalTok{(airquality)) }\SpecialCharTok{\textgreater{}} \DecValTok{0}\NormalTok{, ]}
\end{Highlighting}
\end{Shaded}

\textbf{Answer}

\begin{verbatim}
    Ozone Solar.R Wind Temp Month Day
5      NA      NA 14.3   56     5   5
6      28      NA 14.9   66     5   6
10     NA     194  8.6   69     5  10
11      7      NA  6.9   74     5  11
25     NA      66 16.6   57     5  25
26     NA     266 14.9   58     5  26
27     NA      NA  8.0   57     5  27
32     NA     286  8.6   78     6   1
33     NA     287  9.7   74     6   2
34     NA     242 16.1   67     6   3
35     NA     186  9.2   84     6   4
36     NA     220  8.6   85     6   5
37     NA     264 14.3   79     6   6
39     NA     273  6.9   87     6   8
42     NA     259 10.9   93     6  11
43     NA     250  9.2   92     6  12
45     NA     332 13.8   80     6  14
46     NA     322 11.5   79     6  15
52     NA     150  6.3   77     6  21
53     NA      59  1.7   76     6  22
54     NA      91  4.6   76     6  23
55     NA     250  6.3   76     6  24
56     NA     135  8.0   75     6  25
57     NA     127  8.0   78     6  26
58     NA      47 10.3   73     6  27
59     NA      98 11.5   80     6  28
60     NA      31 14.9   77     6  29
61     NA     138  8.0   83     6  30
65     NA     101 10.9   84     7   4
72     NA     139  8.6   82     7  11
75     NA     291 14.9   91     7  14
83     NA     258  9.7   81     7  22
84     NA     295 11.5   82     7  23
96     78      NA  6.9   86     8   4
97     35      NA  7.4   85     8   5
98     66      NA  4.6   87     8   6
102    NA     222  8.6   92     8  10
103    NA     137 11.5   86     8  11
107    NA      64 11.5   79     8  15
115    NA     255 12.6   75     8  23
119    NA     153  5.7   88     8  27
150    NA     145 13.2   77     9  27
\end{verbatim}

\bookmarksetup{startatroot}

\chapter{Question 2}\label{question-2}

Find mean, standard deviation, minimum, maximum for each of temperature
and ozone level.

\textbf{Data Manipulation}

Extracts the rows in the data without missing values

\begin{Shaded}
\begin{Highlighting}[]
\NormalTok{valid\_rows }\OtherTok{\textless{}{-}}\NormalTok{airquality[}\FunctionTok{rowSums}\NormalTok{(}\FunctionTok{is.na}\NormalTok{(airquality)) }\SpecialCharTok{==} \DecValTok{0}\NormalTok{, ]}
\end{Highlighting}
\end{Shaded}

\section{Mean}\label{mean}

\textbf{Code}

\begin{Shaded}
\begin{Highlighting}[]
\NormalTok{mean\_temp }\OtherTok{\textless{}{-}} \FunctionTok{mean}\NormalTok{(valid\_rows}\SpecialCharTok{$}\NormalTok{Temp)}
\NormalTok{mean\_ozone }\OtherTok{\textless{}{-}} \FunctionTok{mean}\NormalTok{(valid\_rows}\SpecialCharTok{$}\NormalTok{Ozone)}
\end{Highlighting}
\end{Shaded}

\textbf{Answer:} Mean Temperature:

\begin{verbatim}
[1] 77.79279
\end{verbatim}

\textbf{Answer:} Mean Ozone Level:

\begin{verbatim}
[1] 42.0991
\end{verbatim}

\section{Standard Deviation}\label{standard-deviation}

\textbf{Code}

\begin{Shaded}
\begin{Highlighting}[]
\NormalTok{sd\_temp }\OtherTok{\textless{}{-}} \FunctionTok{sd}\NormalTok{(valid\_rows}\SpecialCharTok{$}\NormalTok{Temp)}
\NormalTok{sd\_ozone }\OtherTok{\textless{}{-}} \FunctionTok{sd}\NormalTok{(valid\_rows}\SpecialCharTok{$}\NormalTok{Ozone)}
\end{Highlighting}
\end{Shaded}

\textbf{Answer:} Standard Deviation of Temperature:

\begin{verbatim}
[1] 9.529969
\end{verbatim}

\textbf{Answer:} Standard Deviation of Ozone Level:

\begin{verbatim}
[1] 33.27597
\end{verbatim}

\section{Minimum}\label{minimum}

\textbf{Code}

\begin{Shaded}
\begin{Highlighting}[]
\NormalTok{min\_temp }\OtherTok{\textless{}{-}} \FunctionTok{min}\NormalTok{(valid\_rows}\SpecialCharTok{$}\NormalTok{Temp) }
\NormalTok{min\_ozone }\OtherTok{\textless{}{-}} \FunctionTok{min}\NormalTok{(valid\_rows}\SpecialCharTok{$}\NormalTok{Ozone)}
\end{Highlighting}
\end{Shaded}

\textbf{Answer:} Minimuim Temperature:

\begin{verbatim}
[1] 57
\end{verbatim}

\textbf{Answer:} Minimuim Ozone Level:

\begin{verbatim}
[1] 1
\end{verbatim}

\section{Maximum}\label{maximum}

\textbf{Code}

\begin{Shaded}
\begin{Highlighting}[]
\NormalTok{max\_temp }\OtherTok{\textless{}{-}} \FunctionTok{max}\NormalTok{(valid\_rows}\SpecialCharTok{$}\NormalTok{Temp) }
\NormalTok{max\_ozone }\OtherTok{\textless{}{-}} \FunctionTok{max}\NormalTok{(valid\_rows}\SpecialCharTok{$}\NormalTok{Ozone)}
\end{Highlighting}
\end{Shaded}

\textbf{Answer:} Maximuim Temperature:

\begin{verbatim}
[1] 97
\end{verbatim}

\textbf{Answer:} Maximuim Ozone Level:

\begin{verbatim}
[1] 168
\end{verbatim}

\bookmarksetup{startatroot}

\chapter{Question 3}\label{question-3}

\section{Calculate Beta estimate}\label{calculate-beta-estimate}

Here, Y is the response variable, and X is the design matrix. The cars
data (an R data set, also always available in R) contains two variables:
speed and distance to stop. Fit a simple linear regression model to
these data, i.e.~find the beta estimates, using the equation above, and
matrix calculations in R.

\textbf{Code}

\begin{Shaded}
\begin{Highlighting}[]
\DocumentationTok{\#\#\#\#\#\#\# Create design matrix \#\#\#\#\#\#\# }
\NormalTok{Y }\OtherTok{\textless{}{-}}\NormalTok{cars}\SpecialCharTok{$}\NormalTok{dist}

\NormalTok{intercept}\OtherTok{\textless{}{-}} \FunctionTok{rep}\NormalTok{(}\DecValTok{1}\NormalTok{,}\DecValTok{50}\NormalTok{)}
\NormalTok{X }\OtherTok{\textless{}{-}} \FunctionTok{cbind}\NormalTok{(intercept,cars}\SpecialCharTok{$}\NormalTok{speed)}

\DocumentationTok{\#\#\#\#\#\#\# Calculation \#\#\#\#\#\#\# }
\NormalTok{beta\_hat }\OtherTok{\textless{}{-}} \FunctionTok{solve}\NormalTok{(}\FunctionTok{t}\NormalTok{(X) }\SpecialCharTok{\%*\%}\NormalTok{ X) }\SpecialCharTok{\%*\%} \FunctionTok{t}\NormalTok{(X)}\SpecialCharTok{\%*\%}\NormalTok{ Y}
\end{Highlighting}
\end{Shaded}

\textbf{Answer Beta estimate}

\begin{verbatim}
[1] -17.57909
\end{verbatim}

\bookmarksetup{startatroot}

\chapter{Question 4}\label{question-4}

\section{Confirm Beta estimate by fitting
model}\label{confirm-beta-estimate-by-fitting-model}

Check the same β estimates as when fitting the linear regression model
using lm() in R.

\textbf{Code}

\begin{Shaded}
\begin{Highlighting}[]
\DocumentationTok{\#\#\#\#\#\#\# Create design matrix \#\#\#\#\#\#\# }
\NormalTok{Y }\OtherTok{\textless{}{-}}\NormalTok{cars}\SpecialCharTok{$}\NormalTok{dist}

\NormalTok{intercept}\OtherTok{\textless{}{-}} \FunctionTok{rep}\NormalTok{(}\DecValTok{1}\NormalTok{,}\DecValTok{50}\NormalTok{)}

\NormalTok{X }\OtherTok{\textless{}{-}} \FunctionTok{cbind}\NormalTok{(intercept,cars}\SpecialCharTok{$}\NormalTok{speed)}

\DocumentationTok{\#\#\#\#\#\#\# Fit Model \#\#\#\#\#\#\# }
\NormalTok{model }\OtherTok{\textless{}{-}} \FunctionTok{lm}\NormalTok{(Y }\SpecialCharTok{\textasciitilde{}}\NormalTok{ cars}\SpecialCharTok{$}\NormalTok{speed , }\AttributeTok{data =}\NormalTok{ cars)}
\end{Highlighting}
\end{Shaded}

\textbf{Answer}

\begin{verbatim}

Call:
lm(formula = Y ~ cars$speed, data = cars)

Coefficients:
(Intercept)   cars$speed  
    -17.579        3.932  
\end{verbatim}


\backmatter


\end{document}
